\section{Approximation for the low exponents}\label{sec:low-assembly}

We prove the approximation theorem for $1\le p\le2$.  The local
constructions in this section use the curve and surface budgets of the
preceding section.  Surface budgets, and the inverse-$BV$ theorem, are
used only when $p=2$.

\begin{theorem}[Low-exponent approximation]\label{la:approximation}
Let $\Omega,\Omega'\subset\R^3$ be bounded domains, let $1\le p\le2$, and
let $f:\Omega\to\Omega'$ be a homeomorphism in
$W^{1,p}(\Omega;\R^3)$.  For every $\eps>0$ there is a locally
bi-Lipschitz homeomorphism $H:\Omega\to\Omega'$ such that
\begin{equation}\label{la:goal}
 \int_\Omega \bigl(|H-f|^p+|DH-Df|^p\bigr)<\eps.
\end{equation}
In particular, the target of $H$ is exactly $\Omega'$, and
$H\in W^{1,p}(\Omega;\R^3)$.
\end{theorem}

We work on finitely many compact sets carrying controlled jets.
At rank three, radial blocks recover those jets; at ranks one and two,
kernel compression reduces the core estimate to a line or plane section,
with a further planar radial construction at rank two when $p<2$.
Boundary parametrization completes the remaining cells, and the final
proof records the order of choices and proves global summability.

Throughout the section, all modifications are made inside the open
domains.  Constants denoted by $C$ depend only on the dimension, $p$,
and the fixed shape constants.  A subscript records any additional
dependence.  In an expression $o(1)$ involving a mesh size, all the
other parameters and the finite compact configuration have already
been fixed.  The quantities tending to zero will always be
nonnegative costs or upper bounds for such costs.

\subsection{Composition at a collapsed stratum}

We first record the regularity fact used when a three-dimensional
region is compressed toward a line or a plane.  It also explains why
the finite-piece regularity in the budget construction is retained.

\begin{lemma}[Tangential composition limit]\label{la:composition-limit}
Let $D$ be a compact Lipschitz parameter piece of dimension $d$, and let
$q_j,q_0:D\to U$ be Lipschitz maps into an open subset of a Euclidean
space.  Suppose that their images lie in one compact subset of $U$,
that $q_j\to q_0$ uniformly, and that
\[
 Dq_j\longrightarrow Dq_0\quad\hbox{a.e. on }D,
 \qquad \sup_j\|Dq_j\|_{L^\infty(D)}<\infty.
\]
Let $v:U\to\R^a$ be Lipschitz on a neighborhood of that compact set.
Suppose there are finitely many closed pieces $C_\nu$ covering the
neighborhood and $C^1$ maps $v_\nu$, defined on neighborhoods of
$C_\nu$, such that $v=v_\nu$ on $C_\nu$.  Then, for every finite
$r\ge1$,
\begin{equation}\label{la:composition-convergence}
 D(v\circ q_j)\longrightarrow D(v\circ q_0)
       \quad\hbox{in }L^r(D).
\end{equation}
No rank assumption is imposed on $Dq_0$.  The assertion applies
separately to an edge, a face, or a volume piece.
\end{lemma}

\begin{proof}
For each $j\ge0$ and each $\nu$, apply the density theorem on the
parameter piece to $q_j^{-1}(C_\nu)$.  At almost every point of this
preimage, locality of the weak derivative gives
\[
 D(v\circ q_j)=Dv_\nu(q_j)Dq_j.
\]
Remove the union of the exceptional sets for the countably many $j$
and finitely many $\nu$.  At a remaining parameter point, every branch
containing $q_0$ gives the same product with $Dq_0$, namely
$D(v\circ q_0)$.  If a branch occurs along infinitely many $q_j$, its
closedness implies that it contains $q_0$.  Continuity of its
derivative and convergence of $Dq_j$ therefore show that its products
converge to this common value.  Finiteness of the branch list proves
pointwise convergence of the full sequence.  The fixed derivative
bounds give a bounded majorant, so dominated convergence proves
\eqref{la:composition-convergence}.  The proof takes place on $D$;
it does not pull back an ambient null set through a possibly
rank-deficient map.
\end{proof}

\subsection{Full-column-rank radial blocks}

Here the parameter dimension is either $m=3$, with $1\le p\le2$, or
$m=2$, with $1\le p<2$.  In the latter case the parameter domain is a
flat plane section of the source.  All target motions remain ambient
three-dimensional homeomorphisms.

Let $K$ be a compact subset of such a parameter patch.  Assume that
there is a local $C^1$ ambient diffeomorphism $b$ such that
\begin{equation}\label{la:identity-jets}
 f=b,\qquad D_m f=D_m b\quad\hbox{on }K
\end{equation}
up to the null sets of the indicated traces.  The map $f$ on the patch
is Sobolev and has its continuous ambient values.  Both $Db$ and
$(Db)^{-1}$ are bounded by a fixed constant, denoted collectively by
$C_b$.  For $m=2$, the additional normal column of $b$ is chosen so
that its ambient orientation agrees with that of $f$.
Write
\[
 F=b^{-1}\circ f.
\]
Thus $F=\Id$ and $D_mF=I_m$ on $K$, where $I_m$ is the derivative
of the identity inclusion of the parameter plane into $\R^3$.
Target coordinates in the following construction are always the
$b$-coordinates.  If $m=3$ and $p=2$, then
\begin{equation}\label{la:inverse-bv}
 F^{-1}\in BV_{\mathrm{loc}},
\end{equation}
by \cite[Theorem~1.1]{CsornyeiHenclMaly2010}.  That theorem requires no
finite-distortion hypothesis.

Choose an integer $N$, put $l=N^{-1}$, and fix a sufficiently large
absolute constant $K_0$.  Set
\begin{equation}\label{la:gamma}
 \gamma=l/K_0.
\end{equation}
Use a uniformly shifted cubical grid of side $h$, with each macro cube
subdivided into tiles of side $lh$.  The word ``cube'' includes a
square when $m=2$.  Cubes meeting $K$ are candidates.  In each candidate
$Q$, draw a center $a$ uniformly in a central cube of side $h/4$; retain
only centers belonging to $K$.  Rays from $a$ through boundary mesh
vertices are the forward spokes.  When $m=3,p=2$, also choose,
independently on each boundary mesh square, a point $u_j$ uniformly in
its central half and use the target ray from $a$ through $u_j$.

Cubical polar coordinates are written
\[
 x=a+r(\omega-a),\qquad \omega\in\partial Q,
\]
so the boundary has radial coordinate $r=1$.  In dimension two, the
ambient output box is the normal thickening of $Q$, centered on the
parameter plane and of height $h$.

\begin{lemma}[Radial block construction]\label{la:radial-blocks}
Fix finitely many separated compact configurations satisfying
\eqref{la:identity-jets}, and fix $l$.  For sufficiently small $h$ one
can retain candidate blocks with the following properties.

\begin{enumerate}
\item The expected $m$-volume of discarded candidate blocks tends to
zero as $h\to0$.  The union of retained blocks therefore differs
from $K$ by a set whose expected measure tends to zero.
\item Each retained $Q$ has the core
\[
 B_Q=a+(1-8\gamma)(Q-a).
\]
The complement consists of shape-controlled frusta of scale $lh$, one
over each boundary tile.  After the common budget construction and
the source and target changes described below, the map on $B_Q$ is
$b$, except on caps of total measure at most $C\gamma |Q|$.  Its
tangential derivative on the whole core and its core boundary is
bounded by $C_b$.
\item Every radial connector in the final shell has image length at
most $C_b\gamma h$.  Nonincident mesh edges avoid the radial-angular
expansions.  Subsequent cap contractions multiply their length
budgets by at most a fixed constant.  With $\ell=lh$, the converted
shell edge cost after constant-speed parametrization is
$\ell^{m-p}\sum_e L_e^p$, where $L_e$ is the image length.
Its bound is $O_b(l)\sum_Q|Q|+o(1)$ after the source selection and
individual curve-budget bounds, with arbitrary additional transfer
errors.
\item If $m=3,p=2$, include the full wedge from $a$ to each boundary
tile edge among the measured surfaces.  There are bounded continuous
weights, fixed before the target budget mesh, for which the final
connector-face areas are bounded by the weighted $H_0$-areas, up to
arbitrarily small errors.  After multiplication by $lh$, the expectation of their total converted budget
$\mathcal B_{\mathrm{wedge}}$ satisfies
\begin{equation}\label{la:wedge-cost}
 \mathbb E\mathcal B_{\mathrm{wedge}}
       \le C_b\gamma\sum_i |K_i|+o(1).
\end{equation}
Outer ordinary edges and faces have converted cost
$O_b(l)\sum_Q|Q|+o(1)$ on these configurations.
\end{enumerate}
The construction can be performed simultaneously with the other
budgeted pieces used below.  In a volume block the only measured
interior surfaces are its wedges.  In the plane-section application,
any additional graph to be fixed is assumed to have positive
clearance from $K$ and its image.  Taking $h$ smaller then keeps all
expansion and cap supports off that graph.  More precisely, write
$V$ for the ambient initial map after its output motions and
$q_Q:Q\to Q$ for the radial parameter change.  The controlled trace
is $V\circ q_Q$.  For $m=3$ these parameter changes are ambient
source self-homeomorphisms; for $m=2$ they are retained solely as
section parametrizations.  The graph preservation means that the
radial motions leave its already-budgeted $H_0$-image unchanged;
it does not assert equality of $H_0$ with $f$ on that graph.
\end{lemma}

\begin{proof}
We give the selection, the maps, and the estimates in their order of
use.

\paragraph{Selection of centers and stars.}
Averaging the macro shift and the center draw is equivalent, up to
fixed bounded densities, to integration in the actual center $a$ and
normalized relative offset variables.  The weight in a block sum is
$h^m$.  Failure to choose $a\in K$ on candidate blocks costs at most
\[
 C\bigl|N_{Ch}(K)\setminus K\bigr|,
\]
which tends to zero.  Here $N_t(K)$ denotes the parameter-space
$t$-neighborhood of $K$.

For almost every center in $K$ and almost every fixed choice of
relative variables, each of the finitely many forward line traces is
absolutely continuous and has derivative the identity at $a$ in the
one-dimensional Lebesgue-point sense.  This follows from slicing and
$D_mF=I_m$ on $K$.  In the reverse system, use
\eqref{la:inverse-bv} and $F^{-1}=\Id$ on $K$, together with directional
$BV$ slicing and its separate absolutely continuous and singular
variation identities
\cite[Section~3.11, Theorems~3.103, 3.107, and~3.108]{AmbrosioFuscoPallara2000}.
At almost every line
density point of $K$, the absolutely continuous derivative of the
inverse trace is the identity and the density of its singular
variation is zero.  The trace is continuous because $F^{-1}$ is
continuous.  Consequently both systems have relative conical
differentiability at $a$, and the length of an initial ray interval is
its radius times $1+o(1)$.

For completeness, near-minimal length gives the single-crossing
sections required by Proposition~\ref{lt:star}.  The radial distance along
one distorted initial ray attains every level up to $(1-o(1))t$ and
has total variation at most $(1+o(1))t$.  One-dimensional coarea shows
that levels with more than one crossing have measure $o(t)$.  A finite
union of these exceptional sets still has measure $o(t)$, so there
are arbitrarily small common levels crossed once by every ray in
the system.  Injectivity and the local differentiability estimate
exclude remote tails from these shrinking sections.

The loop condition in Proposition~\ref{lt:star} is also obtained by slicing.
Fix one of a countable library of finite piecewise smooth loop
systems on the angular sphere.  For a fixed library loop $\omega$, set
\[
 E_t(a)=\int\bigl(|D_mF(a+t\omega(s))-I_m|^p
       +\mathbf1_{\text{patch}\setminus K}(a+t\omega(s))\bigr)
                         |\omega'(s)|\,ds.
\]
Extend the integrands to a fixed neighboring patch.  Translation
continuity and their vanishing on $K$ give
$\int_K E_t(a)\,da\to0$.  A diagonal choice of $t_n\downarrow0$
with summable means for the first $n$ library loops gives
$E_{t_n}(a)\to0$ for every loop at almost every center.  Each loop
then has a point in $K$, where $F=\Id$.  Absolute continuity from
that point gives uniform convergence of
$(F(a+t_n\omega)-a)/t_n$ to $\omega$, also when $p=1$.
The ray estimates hold at all sufficiently small scales, so this
subsequence supplies the loop tests simultaneously with them.
Choose the library fine enough to avoid the separated cones and test
the required punctured-sphere generators.  In the planar case
equatorial circles suffice.  Countability permits all these tests at
the same generic centers.

For any prescribed small cone and positional tolerances depending on
$l$, dominated convergence in the center and relative variables now
shows that the failed tests have total block volume tending to zero.
In full dimension the oriented generator test itself forces the
comparison link to have the identity orientation; no separate
Jacobian-sign assertion is needed here.  In the planar case the
orientation was fixed by the normal column of $b$.

\paragraph{Excluding intrusive edges.}
Call the unmodified fine lattice skeleton the default skeleton,
whether or not a candidate block will later be replaced.  Any default
edge whose image can enter a proposed inset target support lies in a
shrinking source neighborhood $U_h=N_{\rho(h)}(K)$ of $K$, by
continuity of the inverse, where $\rho(h)\to0$.  Enlarge $\rho(h)$
by $O(h)$ to include each entire relevant edge.  No estimate
$\rho(h)=O(h)$ is needed or assumed.
For an edge $e$ meeting $K$, absolute continuity from an unchanged
point gives
\[
 \sup_e|F-\Id|>c\gamma h
 \quad\Longrightarrow\quad
 \int_e|D_t(F-\Id)|^p\ge c(l,\gamma,p)h.
\]
Translation averaging makes the sum of these integrals in the
shrinking neighborhood $o(h^{1-m})$ for fixed $l$.  Edges missing $K$
are charged by their entire length $lh$ to
$U_h\setminus K$.  Thus the expected number of bad edges is
$o(h^{-m})$.  Their total image length is $o(h^{1-m})$: use the same
derivative-error estimate and H\"older on the bad edges, and use the
original derivative integral on those wholly off $K$.  For $p=1$
this is the direct length integral, with no H\"older gain required.

A rectifiable curve of length $L$ meets at most $C(1+L/h)$ grid blocks
at scale $h$, by division into subarcs of length at most $h$.  In the
planar case first project the curve to the parameter plane in target
coordinates.  Discarding all blocks struck by bad-edge images
therefore loses $o(1)$ total volume.  Good nonincident edges are
$o(\gamma h)$ from their original positions and avoid the strictly
inset expansion supports.  Accurate spoke positions also prevent a
spoke of one retained block from entering another block's support.

\paragraph{The protected germ and the radial maps.}
Apply Proposition~\ref{lt:star} at the retained centers.  In the
subsequent use of Lemma~\ref{lt:relative-germs} and
Proposition~\ref{lt:budget-transfer}, require the initial map $H_0$
to satisfy
\[
 b^{-1}\circ H_0=\Id\quad\hbox{near every retained center}.
\]
Keep the narrow forward cones and the small positional errors.  In
the critical case, keep also the condition that the preimage of each
chosen target ray, up to $r\le1-\gamma$, lies in the corresponding
open radial sector inside $Q$, except at $a$.  The star modification
preserves the angular conditions near the center.  Away from a
smaller center neighborhood the constraints have compact positive
margins, and fine approximation with reciprocal control preserves
them.  The simplex displacement is turned off still nearer the
germ.  Thus the later budget transfer does not destroy the ray
clearance.  For now fix the germ neighborhoods and margins.  The
actual budget transfer is made only after the outer weight below
has been fixed; the next paragraphs specify the maps that will
then follow it.

Once $H_0$ is chosen, choose $\lambda>0$ so small that $a+\lambda(Q-a)$ is in the exact
germ.  Precompose on $Q$ with the radial homeomorphism fixing
$\partial Q$ and sending $B_Q$ linearly onto this tiny cube.  On the
remaining ray segments use the linear radial interpolation.  In
target $b$-coordinates postcompose with a radial expansion $r\mapsto
e(r)$ which is the identity for $r\ge1-2\gamma$ and satisfies
\[
 e(r)=\frac{1-8\gamma}{\lambda}r\qquad(0\le r\le\lambda).
\]
This gives exactly $b$ on $B_Q$.  For $p<2$, any strictly increasing
piecewise linear profile between these prescribed portions is
sufficient.  The supports for different blocks are disjoint.

\paragraph{Angular attenuation when $p=2$.}
For this paragraph $m=3$.  On a boundary square $S_j$, write its
points in rays from $u_j$ as
\[
 \omega=u_j+\rho R_j(\theta)v(\theta),\qquad 0\le\rho\le1.
\]
Here $R_j(\theta)$ is the distance to the square boundary.  It is
comparable to $lh$ and is piecewise smooth with the corresponding
uniform derivative bounds.  Choose $0<\kappa,t<1/4$ and
replace $\rho$ by
\[
 \phi(\rho)=
 \begin{cases}
  (1-t)\rho/\kappa,&0\le\rho\le\kappa,\\
  1-t+t(\rho-\kappa)/(1-\kappa),&\kappa\le\rho\le1.
 \end{cases}
\]
The resulting map $A$ preserves $S_j$, fixes its boundary, and
expands its tiny central copy.  On each smooth angular piece,
\begin{equation}\label{la:angular-det}
 \det DA=\phi'(\rho)\frac{\phi(\rho)}{\rho},
 \qquad
 |DA|\le C\left(1+\frac{lh}{|\omega-u_j|}\right).
\end{equation}
The angular dependence of $R_j$ adds shear but leaves the displayed
determinant unchanged.  On a compact set missing $u_j$, choose
$\kappa$ below its relative distance and then let $t\downarrow0$;
the determinant tends uniformly to zero there.  For the straight
interpolation $A_\beta=(1-\beta)\Id+\beta A$, the radial profile is
$(1-\beta)\rho+\beta\phi(\rho)$.  Away from the central copy its
radial slope is at most one, while its transverse expansion has the
bound in \eqref{la:angular-det}.  Hence
\begin{equation}\label{la:angular-interpolation}
 |\det DA_\beta|
 \le C\left(1+\frac{lh}{|\omega-u_j|}\right),
 \qquad0\le\beta\le1,\quad\rho\ge\kappa.
\end{equation}

Choose $0<\lambda<\lambda'<1/4$ still inside the exact affine region.
Turn $A$ on between $\lambda$ and $\lambda'$, leave it on through
$1-4\gamma$, and turn it off between $1-4\gamma$ and $1-3\gamma$.
Choose a slope $0<\alpha<\gamma$ on $[\lambda,1/2]$, and put
\[
 k_\alpha=
 \frac{6\gamma-\alpha(1/2-\lambda)}{1/2-2\gamma}>0.
\]
The denominator is bounded below, so $k_\alpha\le C\gamma$.
On $[1/2,1-2\gamma]$ use this slope.  Explicitly, these portions are
\[
 e(r)=
 \begin{cases}
  1-8\gamma+\alpha(r-\lambda),&\lambda\le r\le1/2,\\
  1-8\gamma+\alpha(1/2-\lambda)+k_\alpha(r-1/2),
       &1/2\le r\le1-2\gamma.
 \end{cases}
\]
They meet $e(1-2\gamma)=1-2\gamma$ exactly.  The slope $\alpha$
can still be decreased as required by the inner error estimate.
All maps for fixed positive parameters are bi-Lipschitz with
finitely many closed $C^1$ pieces.

In normalized cubical polar coordinates the target map is
\[
 (r,\omega)\longmapsto
       a+e(r)\bigl(A_{\beta(r)}(\omega)-a\bigr).
\]
Its two-angular-direction area factor is bounded by
\begin{equation}\label{la:two-angular}
 C_b(e/r)^2|\det DA_{\beta(r)}|,
\end{equation}
and its mixed radial-angular factor is bounded by
\begin{equation}\label{la:mixed-angular}
 C_b(e/r)\bigl(|e'|+e\,l|\beta'|\bigr)|DA_{\beta(r)}|.
\end{equation}
Indeed the radial derivative is
$e'(A_\beta-a)+e\beta'(A-\Id)$, the angular derivative is $eDA_\beta$,
and $|A-\Id|\le Clh$.  Uniform transversality of the cubical radial
direction to a boundary tile gives these estimates for arbitrary
tangent two-planes as well.

These two estimates treat different components of a surface tangent
plane.  Its purely angular component is controlled by the determinant
of $DA_{\beta(r)}$, while its mixed component uses only the first power
of $|DA_{\beta(r)}|$.  Write
\[
 G(\omega)=1+\frac{lh}{|\omega-u_j|}
\]
on each boundary tile.  Figure~\ref{la:radial-bands} displays the
resulting area bounds in the unexpanded target coordinates.  The inner
costs will be made small after $H_0$ has been constructed, using its
actual measured surfaces.  The remaining costs will be bounded by a
continuous outer weight fixed before the budget transfer.  This is why
the later attenuation choices do not change the budget already selected.

\begin{figure}[tb]
\centering
\begin{tikzpicture}[x=0.96cm,y=1cm,font=\scriptsize]
 \fill[blue!18] (0,0) rectangle (1,0.85);
 \fill[yellow!20] (1,0) rectangle (2,0.85);
 \fill[blue!5] (2,0) rectangle (5.8,0.85);
 \fill[blue!13] (5.8,0) rectangle (9,0.85);
 \fill[orange!25] (9,0) rectangle (10.6,0.85);
 \fill[orange!12] (10.6,0) rectangle (12.2,0.85);
 \fill[gray!10] (12.2,0) rectangle (14,0.85);
 \draw (0,0) rectangle (14,0.85);
 \foreach \x in {1,2,5.8,9,10.6,12.2} {\draw (\x,0)--(\x,0.85);}
 \node[align=center] at (0.5,0.42) {linear\\core};
 \node[align=center] at (1.5,0.42) {turn\\on $A$};
 \node[align=center] at (3.9,0.42) {arbitrarily small area\\after $H_0$ is fixed};
 \node at (7.4,0.42) {$O(\gamma G)$};
 \node at (9.8,0.42) {$O(G)$};
 \node at (11.4,0.42) {$O(1)$};
 \node at (13.1,0.42) {identity};
 \node[above] at (7.4,0.88) {$A$ fully on};
 \node[above,align=center] at (9.8,0.88) {turn off\\$A$};
 \draw[-{Stealth[length=2mm]}] (0,-0.08)--(14.5,-0.08)
       node[right] {$r$};
 \foreach \x/\s in {0/0,1/\lambda,2/\lambda',5.8/\tfrac12,
                     9/1-4\gamma,10.6/1-3\gamma,12.2/1-2\gamma,14/1}
       {\draw (\x,-0.02)--(\x,-0.15);\node[below] at (\x,-0.17) {$\s$};}
 \draw[decorate,decoration={brace,mirror,amplitude=4pt}]
       (9,-0.78)--(14,-0.78);
 \node[below] at (11.5,-0.88) {outer band of width $O(\gamma)$};
\end{tikzpicture}
\caption{Critical radial area costs in the pre-expansion coordinate
$r$; intervals are not drawn to scale.  Here
$G=1+lh/\dist(\omega,u_j)$ has bounded angular mean.  The linear core
contains no uncontrolled connector face.  Inner costs can be
suppressed after $H_0$ is fixed.  Among the remaining terms, the middle
band carries the factor $\gamma$, and terms without this factor occur
only in the thin outer band.}
\label{la:radial-bands}
\end{figure}

The measured connector surfaces after source squeezing are truncated
wedges.  Between $\lambda$ and $\lambda'$ they are exactly radial
over tile edges and are unchanged by $A$; their cost is arbitrarily
small with the increase of $e$ on that interval.  Below $\lambda$
these surfaces are absent.  On $[\lambda',1/2]$, all relevant
surfaces have positive angular clearance from the selected rays.
After $H_0$ and $\lambda'$ are fixed, choosing $\kappa,t$ sufficiently
small suppresses \eqref{la:two-angular}, and choosing $e'$ sufficiently
small suppresses \eqref{la:mixed-angular}.  Only finiteness of the
fixed surface areas is used in this step.  More explicitly, work in
units $h=1$ and let
\[
 M_0=\int_{\lambda'\le r\le1-4\gamma}G(\omega)\,d\mu<\infty,
\]
where $\mu$ is the unexpanded $H_0$-area of the finitely many surfaces.
Choose $\kappa$ below the angular clearance of every measured
surface on $\lambda'\le r\le1-\gamma$, including the switch-off
band.  This clearance follows from the inverse-ray sector condition
and compactness after removal of the exact germ.  Thus every use of
\eqref{la:angular-interpolation} on these surfaces lies in
$\rho\ge\kappa$.  The suppressed error is bounded by
\[
 C_b\left(\frac{t}{(\lambda')^2}
                +\frac{\alpha}{\lambda'}\right)M_0.
\]
The exact radial wedges on $[\lambda,\lambda']$ cost at most
$C_b L\alpha(\lambda'-\lambda)$, where $L$ is their total angular
edge length.  Thus $t$ and $\alpha$ can make the sum smaller than any
given positive error, while $k_\alpha$ remains positive and bounded
by $C\gamma$.

On the remaining fully-on band the mixed cost is at most
$C_b\gamma(1+lh/|\omega-u_j|)$ times unexpanded area; the
two-angular cost can again be suppressed.  On the switch-off band,
$l|\beta'|\le C l/\gamma=CK_0$, so
\eqref{la:angular-interpolation}--\eqref{la:mixed-angular} bound both
costs by a constant times the \emph{first} power of
$1+lh/|\omega-u_j|$.  The radial-only outer end has a fixed bounded
area factor.  Squaring the angular factor would not give the needed
averaging estimate, and is avoided by the determinant calculation.

\paragraph{A weight fixed before budgeting.}
Inside a retained volume block, the measured surface list consists
only of its full wedges.  Every other measured source surface lies
outside the block interior; kernel templates remain in their own
macro blocks.  The preimage of a selected target ray lies in the
open sector over its tile, so it meets neither these outside
surfaces nor a wedge, which lies on a sector boundary.  On compact
ray segments away from the exact germ, local finiteness therefore
gives positive clearance from the entire measured surface list.

In the outer radial region choose a continuous buffered majorant
$W$.  The angular envelope on a tile is a buffered version of
\[
 1+\frac{lh}{|\omega-u_j|}.
\]
Its singularity can be capped in a tiny neighborhood of the ray:
the ray-preimage sector condition makes the outer compact portions
of all measured surfaces disjoint from that ray.  First choose the
caps for the adjusted topological map with a positive margin, then
choose the budget mesh and $G,T$ fine enough to preserve the margin.
Cutoffs in disks of radius $O(lh)$ and a constant background make
the envelope continuous across tile boundaries.  Write the resulting
bounded envelope as $\widehat G$.  It agrees with an upper bound for
the uncapped factor on the measured surfaces and is bounded above
by a fixed multiple of the uncapped integrable envelope everywhere.

Choose continuous radial cutoffs with the following values, using
linear interpolation between the listed bands:
\[
\begin{array}{c|cc}
 &0&1\\ \hline
 \chi_{\mathrm{low}}&r\le1/2-\gamma&r\ge1/2\\
 \chi_{\mathrm{out}}&r\le1-6\gamma&r\ge1-5\gamma\\
 \chi_{\mathrm{cut}}&r\ge1-\gamma&r\le1-2\gamma.
\end{array}
\]
For a sufficiently large fixed $C_b$, take
\begin{equation}\label{la:outer-weight}
 W=\chi_{\mathrm{out}}+
   C_b\chi_{\mathrm{low}}\chi_{\mathrm{cut}}
             (\gamma+\chi_{\mathrm{out}})\widehat G.
\end{equation}
It is zero for $r\le1/2-\gamma$, dominates the required
$C_b\gamma G$ on the middle band and $C_bG$ on the return ramp,
and equals the ordinary weight one for $r\ge1-\gamma$.
The radial outer end is covered as well, and the angular factors
are confined below $1-\gamma$.  The part without the factor $\gamma$ occupies an
$O(\gamma)$ radial band.  Buffered cutoffs may enlarge these bands
by another fixed multiple of $\gamma$ without changing the estimates.

This fixes a bounded continuous weight before
Proposition~\ref{lt:budget-transfer} is applied.  The weight vanishes in
the star-modification neighborhoods, and the simplex map is the
identity on smaller exact-germ neighborhoods.  Thus the measured
crossings used by the transfer are the original $f$-crossings.  The
mesh-element suprema of $W$ give the weighted $H_0$-area upper bound.
The still smaller attenuation parameters in the preceding paragraph
are chosen \emph{after} $H_0$; they need not have been known when $W$
was fixed.  Equations~\eqref{la:two-angular} and
\eqref{la:mixed-angular} prove the asserted weighted bound for all
connector, outer, or nonincident measured surfaces, with arbitrarily
small additive errors.

\paragraph{Spokes and the final contractions.}
After the expansion, each shell spoke lies in a ball of radius
$C\gamma h$ about its intended outer mesh vertex.  Its entire image is outside the open linear $\lambda$-core by
injectivity and the exact germ, with the inner endpoint on that
core's boundary.  Its angular direction remains in a sufficiently
narrow vertex cone, and along the whole trace its pre-expansion
radius is at most $1+o(\gamma)$.  The angular maps fix vertices and have uniform local
Lipschitz bounds near them.  These facts give the stated ball
inclusion.

Each such spoke has finite length.  Homothetically contract its
entire vertex ball to a sufficiently tiny concentric ball, and extend
radially to the identity on the twice-larger ball.  The extension has
a uniform upper Lipschitz bound, while the homothety factor may be as
small as required to reduce the whole spoke length to
$C_b\gamma h$.  Group all spokes at a shared vertex and use the same
contraction for them.  The twice-larger balls are pairwise disjoint
when $K_0$ is large enough.  Their centers are outside the open cores.
There are $O(l^{1-m})$ boundary vertices, so their intersection with
a core has measure at most
\[
 C l^{1-m}(\gamma h)^m\le C\gamma h^m.
\]
On the core the map before contraction is $b$; hence the contracted
map still has derivative bounded by $C_b$.  All other length and
area estimates are multiplied by a bounded factor.  No uniform
inverse Lipschitz bound in the contraction factor is required.

\paragraph{Averaging the wedge weight.}
The total area of the full wedges in a three-dimensional block is
$O(h^2/l)$.  On their portions in $K$, $F=\Id$.  The outer band has
relative radial thickness $O(\gamma)$, and the remaining nonzero
weight has the factor $\gamma$.  Also, for a point $\omega$ fixed
before the choice of $u_j$,
\[
 \mathbb E_{u_j}\left[1+\frac{lh}{|\omega-u_j|}\right]\le C,
\]
because $u_j$ is sampled in two dimensions at scale $lh$.  An
indicator of successful selection can be dropped in this upper
bound.  Multiplying the resulting weighted area by $lh$ gives
$C_b\gamma h^3$ per block.

On off-$K$ wedge portions, use the same uncapped angular envelope
before making any clearance-dependent choices.  For normalized wedge
parameters $v$, translation continuity gives
\[
 \int_K
  \bigl(|Df|^2\mathbf1_{\Omega\setminus K}\bigr)(a+hv)\,da
       \longrightarrow0,
\]
and integration in the relative variables and wedge parameters
gives the corresponding averaged trace estimate.  The generic
surface formula in Lemma~\ref{lt:trace-slicing} applies.  This proves
\eqref{la:wedge-cost}, including the later target-mesh averaging of
Lemma~\ref{lt:budget-expectation}.  The outer grid has only
$O(l^{1-m})$ edges or faces at scale $lh$ per macro boundary, so its
converted bounded-derivative cost is $O_b(l)|Q|$.  Its off-$K$ error
tends to zero by the same trace averaging.  Nonincident edges avoid
the expansion; nonincident surfaces use the same clearance and
weighted estimate.  All required assertions follow.
\end{proof}

\subsection{Compact jets and kernel blocks}

We next reduce the remaining positive-rank sets to line and plane
sections.  The distinction between invertible, zero, and nonzero
singular derivatives already plays a central role in the planar
construction of Hencl--Pratelli
\cite[Section~1.1 and Section~3]{HenclPratelli2018}.  Here the singular
models are handled by compressing their kernel directions.  A matrix
with a prescribed rank need not be the derivative of any ambient
diffeomorphism; its kernel is used only to choose the source coordinates.

\begin{lemma}[Finite compact jet selection]\label{la:jet-selection}
Given $\varepsilon_0>0$, one can choose $M<\infty$ and finitely many
pairwise disjoint compact sets $K_i\Subset\Omega$, of ranks
$k_i\in\{1,2,3\}$, such that
\begin{equation}\label{la:jet-tail}
 \int_{\Omega\setminus\bigcup_iK_i}|Df|^p<\varepsilon_0.
\end{equation}
On their nonzero singular values there are common bounds $M^{-1}$
and $M$.  At rank three the sets have ambient $C^1$ diffeomorphism
charts $b_i$ agreeing with the values and derivatives of $f$, with
chart bounds $C_M$.  At ranks one and two, for any subsequently chosen
$\delta>0$, the pieces can be refined so that
\begin{equation}\label{la:rank-model}
 |Df-D_i|\le\delta\quad\hbox{on }K_i,
 \qquad \rank D_i=k_i,
\end{equation}
with the same type of singular-value bounds.  The compact pieces can
be enclosed in disjoint source neighborhoods and corresponding
disjoint target neighborhoods with buffers.  Once all finite
multiplicative constants are fixed, the neighborhoods can be chosen
so that
\begin{equation}\label{la:excess-neighborhood}
 \sum_i\int_{U_i\setminus K_i}(1+|Df|^p)
\end{equation}
is as small as prescribed, and the inverse images of the buffered
target neighborhoods lie in the corresponding $U_i$.
\end{lemma}

\begin{proof}
Sobolev maps are approximately differentiable almost everywhere.
Truncate the positive-rank sets where their nonzero singular values
lie in $[M^{-1},M]$, and choose $M$ so large that the lost derivative
integral is small.  Rank-zero points have zero derivative energy and
need not be included.  On each retained part, pass to compact subsets
with continuous jet data and uniform approximate differentiability,
losing an arbitrarily small further derivative integral.

These are $C^1$ Whitney jets after a further compact refinement.  To
see the value condition, take two nearby retained points and compare
their approximate first-order expansions at a common good point in
the overlap of balls of radius comparable to their separation.  The
uniform density losses can be made smaller than a fixed fraction of
the overlap, so such a point exists.  Subtracting the two expansions,
using continuity of the jets, and then sending the differentiation
error to zero gives the Whitney remainder estimate.  The $C^1$
extension theorem~\cite[Part~I, Section~4, Theorem~I]{Whitney1934} supplies the extensions.  At full rank, continuity
of the derivative and the inverse function theorem give local
invertibility.  Injectivity on the compact set, together with
compactness, makes the extension injective on a sufficiently small
neighborhood of the whole compact piece.  Its derivative and inverse
derivative bounds can be kept within $C_M$.

For ranks one and two, cover the bounded rank-$k$ matrix range by
finitely many $\delta$-balls centered at rank-$k$ matrices with
comparable nonzero singular values.  Refine the compact pieces
accordingly, losing an arbitrarily small integral to make the finite
pieces disjoint.  Their number does not enter the local chart
bounds.  Injectivity of $f$ makes their compact images disjoint as
well.  Regularity of the integrable measure $(1+|Df|^p)\,dx$ gives
\eqref{la:excess-neighborhood}; continuity of $f^{-1}$ permits target
buffers whose preimages stay inside these source neighborhoods.
The initial losses can be allocated to give
\eqref{la:jet-tail}.
\end{proof}

Use Lemma~\ref{lt:source-mesh} with mesh sizes at most $lh$, deep
cubical regions around the $K_i$, and a common macro shift within
each such region.  At rank $k=1,2$, choose orthonormal coordinates
$x=(y,z)$ with $z\in\ker D_i$ and $y\in\R^k$.  At rank three use
Lemma~\ref{la:radial-blocks}.  A failed macro block is left as ordinary
fine cells.  Outside the retained macro blocks all cells are ordinary,
including the locally finite cells approaching $\partial\Omega$.

The anisotropic target grids of Lemma~\ref{lt:sharp-grid} are needed
only at rank one and at rank two when $p=2$.  Their long directions
span $\operatorname{range}D_i$, their aspect parameter is $\eta>0$,
and a linear normalization $L_i$ can be chosen so that
\begin{equation}\label{la:normalization}
 L_iD_i(y,z)=(y,0),\qquad
 \|L_i\|+\|L_i^{-1}\|\le C_M.
\end{equation}
The range and its orthogonal complement remain orthogonal in this
normalization.  Elsewhere the target grid can be isotropic.

Consider one rank-$k$ candidate macro cube $Q$ of side $h$.  Let $B$
be its concentric product cube obtained by removing one fine layer
from each side, let $Y$ be its projection onto the $y$-coordinates,
and let $z_0$ be its central kernel coordinate.  Put
\[
 S=Y\times\{z_0\}.
\]
Construct a Lipschitz map
\begin{equation}\label{la:flat-collapse}
 P_0(y,z)=(y,p_y(z)):Q\longrightarrow Q
\end{equation}
which is the identity on $\partial Q$ and collapses $B$ onto $S$.
For $y\in Y$, collapse the inner kernel cube to $z_0$ and extend
radially and linearly to the outer kernel boundary.  Outside $Y$,
blend with the identity using a piecewise smooth $y$-cutoff.  These
formulas may be chosen semialgebraic and piecewise $C^1$ on closed
pieces, with Lipschitz bound $C_l$.  They preserve $y$ exactly, so
\begin{equation}\label{la:collapse-model-identity}
 D_iDP_0=D_i.
\end{equation}
For $0<\lambda\le1$, put
\begin{equation}\label{la:positive-collapse}
 P_\lambda=(1-\lambda)P_0+\lambda\Id.
\end{equation}
For each fixed $y$ the kernel radial slopes are positive.  The map is
triangular in $(y,z)$, fixes the boundary, and is bi-Lipschitz on $Q$.

Budget the shell edges after $P_0$, and also its shell faces when
$p=2$.  Include $S$, its tile edges, and all lower-dimensional
intersections that occur in these templates.  In source and target
units divided by $h$, require that the sum of the $p$-power
tangential derivative errors from the model on these finitely many
pieces is at most $\tau^p$.  On a $P_0$-image this means the error
from $D_iDP_0=D_i$; on $S$ it means the error from its $y$-columns.
The unprojected section is required to carry its Sobolev trace even
when $k=2,p<2$.

\begin{lemma}[Successful kernel blocks]\label{la:kernel-selection}
For fixed $M,l,\tau$, the total expected volume of candidate
rank-one or rank-two blocks failing these trace tests tends to zero
as first $\delta\to0$ and then $h\to0$, with sufficiently small
excess neighborhoods.  In successful blocks the transferred
unparametrized shell data have converted cost
\begin{equation}\label{la:kernel-shell-budget}
 O_M(l)\sum_Q|Q|+\text{an arbitrarily small error}.
\end{equation}
Here converted edge cost means $(lh)^{3-p}$ times image length to the
power $p$, and, when $p=2$, converted face cost means $lh$ times
image area.  Fixed factors depending on $l$ and $\eta$ are allowed
on the trace errors.
\end{lemma}

\begin{proof}
Translation averaging on the fixed product templates, with block
weight $h^3$, bounds the expected sum of normalized trace errors by
\begin{equation}\label{la:trace-test-expectation}
 C_l\sum_i\int_{N_{Ch}(K_i)}|Df-D_i|^p.
\end{equation}
The measured curve strata, and the measured surface strata when $p=2$,
are covered by Lemma~\ref{lt:trace-slicing}. For the additional
unprojected rank-two plane section when $p<2$, ordinary Sobolev slicing
and translation Fubini supply the $W^{1,p}$ trace used in these tests.
This step requires neither an image-area formula nor a two-dimensional
Lusin property. On $K_i$, the integrand is at most
$\delta^p$; off $K_i$ it is bounded by
$C_M(1+|Df|^p)$, whose excess integral is small.  Markov's inequality
therefore makes the failed volume small relative to the fixed
threshold $\tau^p$.

There are $O(l^{-2})$ shell cells of side $lh$.  In grid units, with
the longest side restored after the anisotropic rescaling,
\eqref{la:collapse-model-identity} prevents inflation of the model
term.  The factors from $DP_0$ and from anisotropy multiply only
$Df-D_i$.  Lemma~\ref{lt:curve-probability}, followed by H\"older
on each curve, gives the edge length-power budgets.  At $p=2$,
Lemma~\ref{lt:budget-expectation} and the squared derivative trace
integrals give the face budgets.  A bounded model derivative costs
$C_M(lh)^3$ per shell cell after conversion; multiplying by
$O(l^{-2})$ gives $C_Mlh^3$ per macro block.  For fixed $l,\eta$, the
error contributions are made as small as desired by $\tau$ and the
budget-transfer tolerances.  All relevant images lie deep inside
their prescribed target patches for small $h$, by the source and
target buffers.  This proves \eqref{la:kernel-shell-budget}.
\end{proof}

\subsection{Parametrizing the retained line and plane sections}

For a successful rank-$k$ block we shall find a bi-Lipschitz
self-homeomorphism $\sigma_Y$ of $Y$, preserving its tile complex,
and a tangential map $h_Y$ into the target.  If $V$ denotes the initial
map $H_0$ after all the required output motions, the relation is
\begin{equation}\label{la:section-parametrization}
 h_Y(y)=V(\sigma_Y(y),z_0).
\end{equation}
The aim is
\begin{equation}\label{la:section-goal}
 \sum_{\text{successful }Q}
   h^{3-k}\int_Y |Dh_Y-D_i|_y^p
       \quad\hbox{arbitrarily small}.
\end{equation}
Only the $y$-columns of $D_i$ occur in this formula.  The factor
$h^{3-k}$ is the volume of the kernel fibers, up to fixed constants.

\paragraph{Rank one.}
On each interval tile, reparametrize the transferred curve at
constant speed in the normalization $L_i$.  By
Lemma~\ref{lt:sharp-grid}, its expected length budget, divided by
the tile length, is at most
\begin{equation}\label{la:sharp-length-budget}
 1+C_M\eta+o_{\tau}(1)+o_{\mathrm{transfer}}(1),
\end{equation}
where $l$ and $\eta$ have already been fixed.  Its budget is also at
least $1-o(1)$: the section trace test gives almost correct endpoint
differences, and a sufficiently fine transfer preserves the values
to a still smaller error.  Consequently the total positive excess,
weighted by tile length times $h^2$, has arbitrarily small expectation.

The constant speeds are bounded by $C_M$ simultaneously, using
Lemma~\ref{lt:curve-probability}.  To see that its grid-unit constants
do not deteriorate with $\eta$, split the normalized trace derivative
into the tangent model and the error.  The model has size bounded by
$C_M$ in longest-side-restored grid units.  The error has $p$-mean at
most $C_{\eta,l}\tau$, and is made small after the grid parameters
are fixed.

Here is the elementary sharpness calculation.  On a unit normalized
interval, write $u$ for the constant-speed curve, $s=|u'|$, and
$d=u(1)-u(0)$.  If $e$ is the desired unit direction, then
\begin{equation}\label{la:length-sharpness}
 \int_0^1|u'-e|^2=s^2+1-2e\cdot d.
\end{equation}
Since $s\le C_M$, $|d-e|$ is small, and $s\ge1-o(1)$, the right side
is bounded by $C_M(s-1)_+$ plus an arbitrarily small error.
For $1\le p\le2$, H\"older converts this squared error to a
$p$-error.  Rescaling the tiles and summing with the fiber weights
proves \eqref{la:section-goal}.  In particular the argument works at
$p=1$; it does not appeal to strict convexity of the $L^1$ norm.

\paragraph{Rank two at the critical exponent.}
Now $k=p=2$.  For every planar tile, use the area budget in the
normalization $L_i$.  Lemma~\ref{lt:sharp-grid}, applied to its
oriented tangent $2$-vector and the squared derivative trace error,
gives expected normalized area at most
\begin{equation}\label{la:sharp-area-budget}
 1+C_M\eta+o_\tau(1)+o_{\mathrm{transfer}}(1).
\end{equation}
There is also the lower bound $1-o(1)$ for every sufficiently fine
generic sample.  Indeed the edge tests make the original boundary
trace uniformly close, after translation and normalization, to the
boundary of the identity square.  The transfer preserves that
closeness.  Orthogonal projection of the transferred disk to the
model plane has degree one at every point a small distance inside
the square.  It therefore covers that inner square, and its area is
at least the inner-square area.  Since the transferred area is no
larger than its budget plus a chosen small error, the same lower
bound holds for the budget.  These tiles lie away from all exempt
star centers.

It follows that the weighted sum of the positive area excesses has
arbitrarily small expectation.  To obtain the needed boundary
parametrization, subdivide each common edge into $J$ fixed intervals
and use constant speed on each interval, preserving its endpoints.
Choose $J$ large before the finer trace and transfer tolerances.
Lemma~\ref{lt:curve-probability} bounds the normalized speeds by
$C_M$; the error part is $C_{\eta,l,J}\tau$ and is made small later.
Every reparametrized point stays within its $J$-interval, so the
boundary values converge uniformly to the ideal translated square
as $J\to\infty$ and the finer errors tend to zero.  The boundary
squared derivative integral is uniformly bounded.

Absorb these common edge changes by coning in each tile, then apply
Lemma~\ref{lt:disk}.  Fix a desired derivative tolerance $\zeta>0$.
The sharp assertion of Lemma~\ref{lt:disk} supplies an area-excess tolerance
$\rho>0$ and a boundary-closeness tolerance, using the already fixed
speed bound.  On tiles whose normalized area is at most $1+\rho$,
it gives squared derivative error at most $\zeta$ in tile units.
For the remaining tiles, the sum of their converted tile measures
is bounded by $\rho^{-1}$ times the positive excess.  The crude disk
estimate bounds their energy by a constant times area plus the
bounded boundary energy.  Thus their total converted energy is at
most
\[
 C_M(1+\rho^{-1})\,
       \bigl(\hbox{total converted positive excess}\bigr)
       +o(1).
\]
Make the positive excess small after fixing $\rho$.  This proves
\eqref{la:section-goal} also in the critical rank-two case.  The
finite-piece and edge regularity needed by the disk estimate will
be checked below before the volume gluing.

\paragraph{Rank two below the critical exponent.}
Suppose $k=2$ and $1\le p<2$.  The macro selection has already fixed
finitely many plane sections.  On each $S$, first discard points
where the tangential derivative differs substantially from the
injective $y$-columns of $D_i$.  On the retained set its two singular
values lie between bounds depending only on $M$.  The trace test and
Chebyshev's inequality show that the discarded integral of
$1+|D_yf|^p$ is at most $C_M\tau^p h^2$.  Select compact $C^1$ Lusin
jet sets in the interiors of the original tiles, losing an
arbitrarily small additional integral.  The jets extend to ambient
charts $b$ by adding a normal column with the ambient orientation of
$f$.  Their derivative and inverse derivative bounds are $C_M$.
Thus the leftover part of the section satisfies
\begin{equation}\label{la:section-leftover}
 h^{-2}\int_{\text{leftover}}(1+|D_yf|^p)
       \le C_M\tau^p+\text{an arbitrarily small error}.
\end{equation}
Choose disjoint compact neighborhoods on the source and target sides.
The new compacts avoid the original tile edges.  They also avoid
$P_0$ of every shell edge: when those images lie in $S$, their
$y$-coordinates belong to the tile boundaries.  Consequently all
these fixed graph images have positive clearance from the new
compact neighborhoods.

Inside each section use a new macro grid of side $s$ and a fine grid
of side $ls$, with $s\to0$ only after the macro data have been fixed.
To respect the original tile seams, insert their fixed coordinate
levels, omit shifted $s$-levels within $s/4$ of a seam, and subdivide
each remaining gap into $N$ equal pieces.  The resulting rectangles
have controlled aspect ratios.  Near a seam, every nonfixed
subdivision level has a shift coefficient of magnitude at least
$1/N$, so its translation density is bounded for fixed $l$.
Only boundedly many such levels occur per seam in an $O(s)$ band.
Therefore their edge trace integrals, multiplied by $ls$, are bounded
in expectation by $C_l$ times the section derivative integral on
that shrinking band.  The fixed seam lines have finite trace
integrals and their converted weights tend to zero.  Both
contributions vanish as $s\to0$.

Apply Lemma~\ref{la:radial-blocks} with $m=2$ on the new compacts,
simultaneously with the preparation of $H_0$.  All target radial and
vertex motions avoid the already-budgeted fixed graph because of the
clearance just established.  On each successful micro block, make
the radial source change only as a parametrization of the section,
and keep its controlled core.  Parametrize the remaining planar
edges at constant speed, except for the prescribed core edges, and
apply Lemma~\ref{lt:collapse} to the remaining two-dimensional cells.
It applies because $p<2$.  Do the same on fallback micro cells.
These changes fit together to form a tile-preserving bi-Lipschitz
map $\sigma_Y$ of the section.  They are not inserted as an
additional ambient precomposition.

On the compact jets in the micro cores, the map equals $b$, apart
from the bounded-derivative caps of Lemma~\ref{la:radial-blocks}; its
derivative is $D_yf$ there.  The cap fraction and shell contributions
are $O_M(l)$ in section units.  The derivative is bounded by $C_M$
on the remaining core portions.  The shell spokes have image length
$O_M(\gamma s)$, so constant-speed parametrization gives their
required bounded energy.  Ordinary outer and fallback edges use the
individual high-probability length bounds, followed by the planar
collapse estimate.  The latter has the form
\[
 \int_D|Dg|^p
       \le C(ls)\int_{\partial D}|D_tg|^p+\text{a chosen error}.
\]
Translation averaging and \eqref{la:section-leftover} therefore give
\begin{equation}\label{la:micro-section-error}
 h^{-2}\int_Y|Dh_Y-D_i|_y^p
       \le C_M l+C_M\tau^p+o_s(1)
                 +\text{an arbitrarily small error}.
\end{equation}
Here failed micro blocks contribute vanishing cost by their failed
area and bounded compact jets; elsewhere the ordinary converted
trace cost is bounded by the leftover integral.  The constants in
that latter bound depend on $M$, not on later micro-chart widths or
their number.  The seam errors were handled separately above.
Multiplying by $h$ and summing proves \eqref{la:section-goal}.
On the original tile edges, the same construction gives an upper
$p$-energy bound by a fixed constant times the original trace
$p$-energy plus chosen errors.  There are no micro output motions
on those edges.

\subsection{From section cores to volume cells}

We now keep all finite section data and output motions fixed.  On a
rank-$k$ macro core use
\begin{equation}\label{la:kernel-core-map}
 G_\lambda(y,z)
   =V\bigl(P_\lambda(\sigma_Y(y),z)\bigr),\qquad(y,z)\in B.
\end{equation}
Since $P_0$ is flat on $B$, its argument here is
\[
 (\sigma_Y(y),z_0+\lambda(z-z_0)).
\]
Lemma~\ref{la:composition-limit} gives
\begin{equation}\label{la:kernel-core-limit}
 DG_\lambda\longrightarrow[Dh_Y,0]
       \quad\hbox{in }L^p(B)\qquad(\lambda\downarrow0).
\end{equation}
The outer map $V$ has finitely many closed-piece $C^1$ formulas on
the compact configurations: this is true for $H_0$ and for each
radial, angular, and cap map.  The inner maps are uniformly Lipschitz
for fixed $\sigma_Y$, and their derivatives have the stated pointwise
limit.  This verifies the hypotheses even if the limiting map has
rank one or two.

The core boundary values preserve the images of its individual facets.
Indeed, set
\[
 w_\lambda=V\circ P_\lambda,\qquad
 T_B(y,z)=(\sigma_Y(y),z).
\]
Then $G_\lambda=w_\lambda\circ T_B$ on $B$.  Since $\sigma_Y$ preserves
the tile complex, $T_B$ preserves each fine facet and edge of
$\partial B$.  In particular,
\[
 G_\lambda(F)=w_\lambda(F)
 \quad\hbox{for every such facet }F.
\]
Thus these boundary values can be prescribed on the adjacent shell
cells without changing their base face images.  The product formula
need not be extended through the shell; its base map there is
$w_\lambda=V\circ P_\lambda$.
Apply Lemma~\ref{la:composition-limit} separately on its edges and
faces, with limiting parameter map $P_0$.  Their lengths and areas
converge to the transferred budgets, or to smaller upper bounds
when rank is lost.

We record the boundary scaling of the prescribed core data.  Facets
normal to a kernel coordinate have limiting energy bounded by
\begin{equation}\label{la:normal-kernel-face}
 C h^{2-k}\int_Y|Dh_Y|^p.
\end{equation}
For $k=2$, facets normal to a $y$-coordinate use the original tile-edge
parametrizations on $\partial Y$; their total energy per macro block
is $O_M(h^2)$.  For $k=1$ their limiting energy is zero.  The total
limiting edge energy on the subdivided $\partial B$ is
\begin{equation}\label{la:core-edge-energy}
 O_M(h/l).
\end{equation}
For the critical rank-two case this follows from the uniform edge
speed bounds.  Below the critical exponent use instead the
individual original tile-edge energy estimates just proved; the
trace tests bound their sum.  Edges replicated in kernel directions
have the same tangential data, and kernel-axis edges have zero
limiting energy.  Multiplying \eqref{la:normal-kernel-face} by $lh$
and \eqref{la:core-edge-energy} by $(lh)^2$ gives
$O_M(l)$ times the macro-volume contributions, using
\eqref{la:section-goal} for the first expression.  Small positive
$\lambda$ preserves these estimates with arbitrarily small error.

\begin{lemma}[Completion of uncontrolled cells]\label{la:cell-completion}
Let $D$ be an uncontrolled source three-cell of scale $H$ in the
construction.  Its base embedding $v:D\to\R^3$ is bi-Lipschitz and is
given by finitely many $C^1$ formulas extending to neighborhoods of
their closed pieces.  Write $g$ for the prescribed or subsequently
chosen boundary data.  On every prescribed core face $F$,
assume that $g|_F=v|_F\circ s_F$, where $s_F:F\to F$ is a
bi-Lipschitz self-homeomorphism preserving each edge.  These self-maps
agree on shared edges.  When $p=2$, their edge restrictions have
derivatives with essential limits almost everywhere.
Then the cell can be reparametrized bi-Lipschitzly, preserving its
base image and all prescribed boundary data, with derivative cost
at most a constant times
\begin{equation}\label{la:cell-cost}
 \begin{split}
 H\sum_{F\text{ prescribed}}\int_F|D_tg|^p
 &+H^2\sum_{e\subset\partial D}\int_e|D_tg|^p\\
 &+\mathbf1_{\{p=2\}}\,
       H\sum_{F\text{ uncontrolled}}\operatorname{Area}(v|_F)
 \end{split}
\end{equation}
plus any prescribed positive error.  An unprescribed edge of length
comparable to $H$ may be given energy at most
$C\operatorname{length}(v(e))^pH^{1-p}$.
The constants depend only on $p$ and the fixed shape bounds, and
adjacent cells can be completed with identical shared-face data.
\end{lemma}

\begin{proof}
First reparametrize every unprescribed edge at constant speed,
keeping common vertices.  On an unprescribed face, cone-extend these
edge changes.  If $p<2$, Lemma~\ref{lt:collapse} in dimension two gives
face energy bounded by $CH$ times its boundary $p$-energy plus a
chosen error.  If $p=2$, use the crude estimate of
Lemma~\ref{lt:disk}; after scaling it bounds the face energy by a
constant times its base image area plus $H$ times its edge squared
energy.  Shape-controlled polygons are converted to square
parameters by fixed bi-Lipschitz closed-piece smooth charts.
Prescribed faces are left with their given parametrizations.  Make
one choice per shared face.

The disk hypotheses follow from the actual base and edge regularity.
Restricting finitely many closed-piece $C^1$ formulas gives the
almost-everywhere essential limits on a base face and the tangential
limits at its parameter edges.  On each fixed edge subdivision,
normalized arclength has derivative with essential limits almost
everywhere.  Bi-Lipschitzness bounds that derivative away from zero,
and the inverse arclength map has the same property.

For prescribed kernel edges, the source self-map is the restriction
of $T_B$.  In the rank-one case it is the interval arclength change;
in the critical rank-two case it is the arclength change on each fixed
$J$-subinterval of the section edge.  The disk parametrizations fix the
tile boundaries, so they leave these edge changes intact.  Kernel-axis
edges use the identity.  Reusing these source self-maps at positive
$\lambda$ retains their essential derivative limits, even though the
resulting curves need not have constant speed for the base map
$w_\lambda$.  On radial-core facets the prescribed map already equals
$v$, so $s_F=\Id$.  Thus all prescribed edges needed for a critical
face satisfy the stated condition.  Coning the edge changes preserves the
compatibility required by Lemma~\ref{lt:disk}.

The assembled boundary map $g$ gives the face-preserving bi-Lipschitz
self-homeomorphism $v^{-1}\circ g$ of $\partial D$.  Apply
Lemma~\ref{lt:collapse} to this boundary self-map in dimension three,
where $p\le2<3$.
It contributes a factor $H$ times the true boundary-face energy.
Substituting the preceding face estimates gives
\eqref{la:cell-cost}.  This last collapse requires the radial
compatibility of the base parametrization and a fixed bi-Lipschitz
boundary map; it does not require stronger smoothness of the entire
boundary reparametrization.  Allocate the intermediate errors so
their converted sum is below the prescribed cell error.
\end{proof}

\subsection{Global choices, summability, and strong convergence}

\begin{proof}[Proof of Theorem~\ref{la:approximation}]
We first construct locally bi-Lipschitz approximants; smoothing is
performed only after their global derivative costs have been
estimated.  If necessary, compose with a reflection to fix an
orientation convention, and undo it at the end.

\paragraph{Order of the fixed data.}
Let $\varepsilon_0>0$ be a tail tolerance, to be sent to zero.
Choose the compact rank truncation and its bound $M$ in
Lemma~\ref{la:jet-selection}.  Choose $l$ so small that all the
$O_M(l)$ contributions in volume and section shells and in cap
regions are below a prescribed auxiliary tolerance.  Recall that
$\gamma=l/K_0$.

With $M,l$ fixed, choose the sharp derivative tolerances for the
line and plane tiles.  The speed bounds entering those tolerances
are fixed in terms of $M$.  When $k=p=2$, choose the further edge
subdivision $J$ large enough for the desired boundary closeness.
Choose $\eta$ small enough for the sharp length and area coefficients,
then $\tau$ small enough for all trace errors, including the factors
depending on $l,\eta,J$.  Next choose $\delta$ in
\eqref{la:rank-model} small enough, refine the compact pieces, and
choose their nested excess neighborhoods so that
\eqref{la:excess-neighborhood} is small after multiplication by every
now fixed inflation factor.  These bounds do not depend on the
widths of the neighborhoods or on the number of their compact
pieces.  Only then choose $h$ small and make the macro success
selections.  The planar micro configurations, when present, are
chosen afterward, using \eqref{la:section-leftover} and then small
$s$.  Their ambient chart constants are controlled by $M$.

Fix the bounded radial area weights and compact clearance margins
before choosing the target budget mesh.  Choose that mesh using
Proposition~\ref{lt:budget-transfer} and Lemmas~\ref{lt:budget-expectation},
\ref{lt:curve-probability}, and~\ref{lt:sharp-grid}.  This determines
$H_0$, with the required exact affine germs.  The small radial
parameters, angular attenuation parameters, and vertex contractions
are then chosen from the actual finite configuration.  In kernel
blocks choose $\lambda$ after the section parametrizations and all
output maps have been fixed.  Finally complete the uncontrolled
cells using Lemma~\ref{la:cell-completion}, with summable positive
error prescriptions.  Lemma~\ref{pre:local-tolerances} permits the
locally finite fine choices and these summable prescriptions.

\paragraph{Controlling ordinary and failed cells.}
We give the accounting that makes this order of choices sufficient.
For ordinary cells meeting no compact jet set and no enlarged special
neighborhood, source sampling bounds the converted trace integrals in
expectation by
\begin{equation}\label{la:ordinary-tail}
 C_p\int_{\Omega\setminus\bigcup_iK_i}|Df|^p.
\end{equation}
On the simultaneous individual curve-bound event, H\"older's inequality
bounds the converted edge costs by these source trace integrals.  When
$p=2$, the conditional expectation over the later target mesh bounds
the converted face-area budgets by the corresponding source traces.
Substituting in \eqref{la:cell-cost} gives precisely the codimension
conversion factors $H$ for faces and $H^2$ for edges.  The constants in
\eqref{la:ordinary-tail} are ordinary mesh and budget constants,
independent of $M,l,\eta$.

In a failed macro block, the compact-jet integrands are bounded by a
constant depending on $M$.  For the ordinary fine grid, the total
converted size of edges and faces in that block is $O(h^3)$, with
fixed factors from any special target normalization allowed.
The converted source trace integrals over its compact-jet portions
are therefore bounded deterministically by $C h^3$.  After the source
mesh and failed blocks have been selected, the target area-budget
expectation applies to these fixed traces.  On the individual
curve-bound event, H\"older's inequality gives the same bound for
converted edge costs.  Thus the small failed-block volume controls
their total without requiring a source trace law conditional on
failure.  On off-compact portions, discard the failure indicator
before applying the unconditional source trace bounds.

Whenever a nonuniversal constant can occur, its whole source trace
lies in one of the chosen enlarged source neighborhoods for small
enough meshes.  Indeed its image meets a buffered anisotropic or
radial-motion zone, whose inverse image is compactly inside that
neighborhood; the small cell diameter and the fine initial-map
errors preserve this inclusion.  Thus even an edge length-power
estimate can pay its finite factor on the \emph{whole} edge's
H\"older integral, then split this integral into the bounded
compact-jet part and the small excess part.  By
\eqref{la:excess-neighborhood}, all these excess contributions are
as small as prescribed.  No later large constant multiplies the
unprotected tail in \eqref{la:ordinary-tail}.

For full-rank radial blocks, Lemma~\ref{la:radial-blocks} gives the
small wedge budget and the contracted spoke lengths.  On compact
points of a failed full-rank block, $F=\Id$, so their image points
cannot belong to another block's inset angular support.  Outer
ordinary traces of retained blocks have the $O_M(l)$ converted cost
already computed.  Any off-compact surface portion affected by an
angular map is estimated by the uncapped envelope, integrating in
$u_j$ before the later clearance-dependent caps.  Its conditional
mean is bounded for the fixed trace template.  This gives the same
small excess estimate.  Nonincident edges avoid the angular
supports and hence never pay a power of that angular envelope.
Target-mesh suprema of the fixed bounded weights are allowed an
arbitrarily small enlargement error by a still finer target mesh.

The kernel shells are controlled by
Lemma~\ref{la:kernel-selection},
\eqref{la:normal-kernel-face}, and
\eqref{la:core-edge-energy}.  Their prescribed face and edge
contributions in \eqref{la:cell-cost} are $O_M(l)$ in bulk units.
The radial volume shells have the same bound, since their prescribed
core faces have bounded derivatives and their radial edges have
length $O_M(\gamma h)$.  All estimates are relative to the original
straight product or frustum cell sizes; no shape bound is imposed
on the intermediary image cells.

\paragraph{Simultaneous choices.}
Each estimate just used bounds a sum of nonnegative errors.  In the
sharp line and plane cases, this sum is the positive excess, not the
signed total excess alone; the lower bounds in the section argument
are what allow that replacement.  Choose the expectations smaller
than the desired tolerances with enough slack, and use Markov's
inequality for the finitely many aggregate objectives.  The source
and micro configurations are fixed first, using the trace,
failed-volume, and angular-weight bounds.  The target sample is
chosen next, meeting the expectation objectives together with the
arbitrarily high-probability individual curve bounds.  Countably
many ordinary curves are handled with summable failure allowances
and local mesh upper bounds.  Generic slicing and intersection
conditions have full measure on their templates; the star conditions
have already been accounted for by the small failed volume.

In particular, the sharp disk tolerances are selected using fixed
speed bounds.  Make the positive excess small only after their area
threshold $\rho$ is chosen.  The converted measure of bad sharp
tiles and their crude energy are then small by the explicit
$\rho^{-1}$ estimate above.  This avoids any assertion that the
lattice intersection count itself converges pointwise on a
wrinkled trace.  These choices yield deterministic configurations
with all the stated bounds simultaneously.

\paragraph{The resulting homeomorphism.}
Let $V$ be $H_0$ after the ambient radial-angular expansions and the
vertex contractions.  Their supports are disjoint within each
family.  Distinct compact configurations have separated target
neighborhoods; the much later micro-section supports were chosen
smaller still, and avoid all fixed graph data.  In the rank-three
volume blocks precompose by the radial source squeezes.  In kernel
blocks use $P_\lambda$ as the shell base map and
\eqref{la:kernel-core-map} on the core.  Section micro squeezes enter
only through $\sigma_Y$.  Complete all remaining edges, shared faces,
and volume cells as in Lemma~\ref{la:cell-completion}.

Each change is a self-parametrization relative to the same base
image, with matching boundary values.  The meshes are locally finite,
and each interior compact meets only finitely many pieces.  Finite
local straight cell geometry and the finite local bi-Lipschitz
constants give a locally bi-Lipschitz map across the interfaces.
It is a homeomorphism of $\Omega$ onto the same target as $V$.
The initial $H_0$ is onto $\Omega'$, and every output change is a
self-homeomorphism supported inside $\Omega'$.  Consequently the
assembled map $H$ is onto exactly $\Omega'$.

\paragraph{Derivative error.}
Let $\mathcal C$ be the union of the controlled volume cores.  The
preceding bounds imply
\begin{equation}\label{la:outside-core-energy}
 \int_{\Omega\setminus\mathcal C}|DH|^p
       \le C_p\varepsilon_0+o(1),
\end{equation}
where the auxiliary errors, including $O_M(l)$, can be made as small
as prescribed in their stated order.  The original derivative has
the same small bound on this complement: outside the compacts use
\eqref{la:jet-tail}, while on the compacts the shell fraction and
failed volume are small and the jets are bounded.

On a full-rank core, $H=b_i$ except on the small bounded-derivative
caps.  On its compact jet set, $Db_i=Df$; on the remaining portion the
chart derivative is bounded and the excess integral is small.
Thus the derivative difference there is small.  On a kernel core,
\eqref{la:kernel-core-limit} and \eqref{la:section-goal} give closeness
to $D_i$ with the correct fiber factor $h^{3-k}$.  Equation
\eqref{la:rank-model} compares $D_i$ with $Df$ on $K_i$, and the
excess-neighborhood bound treats the remaining points.  Choose the
finitely many $\lambda$'s sufficiently small after the other data.
Combining these estimates with \eqref{la:outside-core-energy} yields
\begin{equation}\label{la:derivative-final}
 \int_\Omega|DH-Df|^p\le C_p\varepsilon_0+o(1).
\end{equation}
All uncontrolled-cell errors were chosen summably.  The displayed
bounds therefore also establish $\int_\Omega|DH|^p<\infty$, rather
than only local Sobolev regularity.

\paragraph{Value convergence.}
The source self-parametrizations move points only within cells or
blocks whose maximum diameter tends to zero.  The initial budget map
can have arbitrarily small value error, and the target motions are
supported in boxes of arbitrarily small diameter in physical
coordinates.  On every interior compact, continuity of $f$ then
gives value convergence.  Since the domains are bounded and all
images lie in the fixed bounded target, dominated convergence gives
$L^p$ convergence on $\Omega$.  Choose the meshes and value
prescriptions so its error and \eqref{la:derivative-final} sum to less
than the prescribed $\eps$; this proves \eqref{la:goal}.

\end{proof}
